% 在过去的文献中，更为常见的连续值序列蒸馏损失函数是沿时序方向进行计算的，目标是最小化每一时刻学生模型输出特征与教师模型输出特征之间差别。如\citep{chang2022distilhubert}中提出的损失函数，最大化同一时刻学生模型输出表示和教师模型输出表示之间的余弦相似度，最小化两者的L1距离来使学生模型学到教师模型中的知识。将这一公式改造成适合我们任务的形式，公式为$$。我们将这种损失函数称之为temporal-axis的，将我们在\ref{sec:032_distillation}中所提出的损失函数称之为dimension-axis的，以此来对两种损失函数进行区分。
In extant literature, the commonly employed loss functions for continuous sequence distillation are typically computed along the temporal axis, with the objective of minimizing the difference between the student and teacher model outputs at each timestep. For instance, the loss function proposed in \citep{chang2022distilhubert} aims to maximize the cosine similarity between the student and teacher model representations at the same timestep while minimizing their $L1$ distance, thereby facilitating the transfer of knowledge from the teacher to the student model. To adapt this formula for our specific task, we can modify the the loss function as follows: 
\begin{align*}
\mathcal{L}_{distll}&=\mathcal{L}_{l1}+\lambda\mathcal{L}_{cos}\\
&=\sum_{t=1}^T[\frac{1}{D}\lVert \mathbf{s}^t-\mathbf{Aq}_1^t\lVert_1 - \lambda\log\sigma(\cos{(\mathbf{s}^t, \mathbf{Aq}_1^t)})],
\end{align*}
where $\mathbf{q}_1^t$ and $\mathbf{s}^t$ respectively denote the quantized output of RVQ first layer and the $D$ dimensional semantic teacher representation at timestep $t$. $\cos(\cdot)$ is cosine similarity. $T$ denotes the number of timesteps and $A$ is the projection matrix. $\sigma(\cdot)$ denotes sigmoid activation. $\lambda>0$ controls the contribution of the cosine layers. We refer to this loss function as "T-axis" to distinguish it from the "D-axis" loss function that we propose in Section \ref{sec:032_distillation}. The latter term is used to denote the loss function introduced in the aforementioned section. These designations are employed to differentiate between these two types of loss functions in this paper.
% 我们探索了两种不同的连续值蒸馏损失函数对\method在\bench上性能的影响，主要结果如表\ref{tab:loss_ablation}所示。相比于表\ref{tab:bench_results}中呈现EnCodec在\bench上的性能，采用"T-axis"的连续值蒸馏损失函数同样使\method在对齐文本的能力上同样获得了很大的提高，但这一能力弱于采用"D-axis"损失函数的\method。在Information Preservation方面，采用"D-axis"损失函数的\method同样强于用用"T-axis"的。实验结果证明"D-axis"连续值蒸馏损失函数取得了比传统的"T-axis"连续值蒸馏损失函数更好的蒸馏效果，我们认为这是由于"D-axis"连续值蒸馏损失函数通过在每个维度上计算余弦相似度，确保学生模型在每个特征维度上都尽可能地靠近教师模型，提供了更多的监督信号促进学生模型的学习，使学生模型不仅关注整体输出的相似性，还关注每个维度的相似性。


We investigated the impact of two distinct continuous distillation loss functions on the performance of \method on \bench. The results of this experiment are summarized in Table \ref{tab:loss_ablation}. When compared to the performance of EnCodec on \bench, as presented in Table \ref{tab:bench_results}, employing the "T-axis" continuous distillation loss function significantly enhances \method's capability in text alignment. However, this improvement is somewhat inferior to that achieved by \method utilizing the "D-axis" loss function. In terms of Information Preservation, \method with the "D-axis" loss function also outperforms its "T-axis" counterpart. The experimental results demonstrate that the "D-axis" continuous distillation loss function yields superior distillation effects compared to the traditional "T-axis" loss function. We attribute this improvement to the "D-axis" loss function's strategy of calculating cosine similarity across each dimension, ensuring that the student model closely aligns with the teacher model on each feature dimension. This approach provides a richer supervision signal, promoting the learning process of the student model by focusing not only on the overall output similarity but also on the similarity within each dimension.

\begin{table*}[h]
    % \setlength{\abovecaptionskip}{-0.cm}
    \setlength{\belowcaptionskip}{-0.2cm}
    % \renewcommand{\arraystretch}{1.2}
    \Large
    \centering
    \resizebox{0.8\textwidth}{!}{
    \begin{tabular}{lc|cc|cc}
        \toprule
        ~ &  & \multicolumn{2}{c}{\textbf{Text Alignment}} & \multicolumn{2}{c}{\textbf{Information Preservation}}   \\
         $\mathcal{L}_{distill}$&  & MI$\uparrow$ & WER$^{\dagger}\downarrow$ & WER$^{\ast}\downarrow$ & SIM$\uparrow$\\
        \midrule
        T-Axis  &RVQ-1 & 26.65 & 21.10  & 10.75 & 0.76 \\
          &RVQ-1:8&25.97  & 21.54  & 5.29 & 0.96 \\
        D-Axis &RVQ-1 & 30.9 & 15.58 & 9.57&0.74  \\
         &RVQ-1:8& 29.7 & 16.03 & 5.04 & 0.97 \\
        
        
        \bottomrule
    \end{tabular}}
    \caption{Results of continuous distillation loss ablation experiment on \bench. We employ the average representation across all HuBERT layers as semantic teachers in this experiment.
    }
    \label{tab:loss_ablation}
\vspace{-0.5em}
\end{table*}
